\section{PSM Placement and Reachability Analysis}\label{sec:reach}

The dVRK patient side manipulator (PSM) holds its tool on a long shaft that
always pivots about one fixed point in space, the \emph{remote center of
motion} (RCM). In surgery this point sits at the incision, so the robot never
pushes sideways on the patient. The drawback is that the tool can only reach a
target by coming in along a line through the RCM, so where the RCM is placed
decides which targets the robot can reach, and from which directions.

Before running anything in simulation, the group used the bracket's known
geometry and the published dVRK kinematics~\cite{kazanzides2014dvrk,sawIntuitive}
to find a good RCM placement. The goal was for the jaw tip to reach
(i)~every point along the inside vertical corner of the L-bracket and
(ii)~every hole in both plates \emph{from the inside}, with the jaws lined up
with the hole. The method has three parts: a model of the arm that can be
solved for the joint angles needed to reach a target
(Sections~\ref{sec:kin}--\ref{sec:ik}), a score that rates a candidate
placement (Section~\ref{sec:score}), and a search that tries many placements
and keeps the best (Section~\ref{sec:search}).

\subsection{Setup and Targets}\label{sec:targets}
All positions are given in the bracket frame $\{B\}$, with $z$ pointing up
and the table at $z=0$. A placement of the RCM is described by six numbers,
\begin{equation}
  \mathbf x=(\mathbf t,\psi,\theta_y,\theta_x),
\end{equation}
where $\mathbf t$ is the RCM's position and the three angles rotate the arm:
$\mathbf R(\mathbf x)=\mathbf R_z(\psi)\mathbf R_y(\theta_y)\mathbf R_x(\theta_x)$.
With all angles zero, the tool points straight down.

Each target $k$ is a point $\mathbf p_k$ that the jaw tip must reach and a
unit direction $\mathbf a_k$ that the jaws must point along. The bracket is
two \SI{4}{mm} thick plates, \SI{65}{mm} tall and \SI{70}{mm} long, and the
positions of all of its features are known from the CAD model:
\begin{itemize}
  \item \textit{Holes} (421, diameter \SI{2.7}{mm}): a $15\times15$ grid on
    one plate and a $14\times14$ grid on the other. The target is the
    centre of the hole on the inside face, and the direction is straight
    through the plate, outward.
  \item \textit{Inside corner} (13 points): spaced every \SI{5}{mm} up the
    inside corner edge, approached diagonally into the corner.
\end{itemize}
For collision checks, each plate is treated as a solid box. The holes are much
smaller than the tool shaft, so this makes no practical difference.

\subsection{Arm Model}\label{sec:kin}
The PSM has six joints,
$\mathbf q=(q_1,\dots,q_6)$: yaw and pitch swing the shaft about the RCM,
insertion $q_3$ slides it in and out, roll $q_4$ spins it, and the wrist
pitch and yaw ($q_5,q_6$) bend the jaws at the end. The jaws are treated as
closed. The arm's dimensions come from the dVRK Denavit--Hartenberg (DH)
table~\cite{craig2005,sawIntuitive} in Table~\ref{tab:dh}. Each row describes
how one joint is placed relative to the previous one, as a fixed transform
\begin{equation}
  {}^{i-1}\mathbf T_i=\mathbf R_x(\alpha_{i-1})\,\mathbf D_x(a_{i-1})\,
  \mathbf R_z(\theta_i)\,\mathbf D_z(d_i),
\end{equation}
where $\mathbf R$ is a rotation and $\mathbf D$ a translation along the named
axis.

\begin{table}[t]
\centering
\caption{dVRK PSM DH parameters and joint ranges~\cite{sawIntuitive}.
Lengths in m.}
\label{tab:dh}
\footnotesize
\setlength{\tabcolsep}{3.5pt}
\begin{tabular}{@{}clcccc c@{}}
\hline
$i$ & Joint & $\alpha_{i-1}$ & $a_{i-1}$ & $\theta_i$ & $d_i$ & Range\\
\hline
1 & Yaw        & $\phantom{-}\pi/2$ & 0      & $q_1+\pi/2$ & 0               & $\pm91^\circ$\\
2 & Pitch      & $-\pi/2$ & 0      & $q_2-\pi/2$ & 0               & $\pm53^\circ$\\
3 & Insertion  & $\phantom{-}\pi/2$ & 0      & 0           & $q_3-0.4318$    & $[0,0.24]$\\
4 & Roll       & 0        & 0      & $q_4$       & $0.4162$        & $\pm260^\circ$\\
5 & Wrist pitch& $-\pi/2$ & 0      & $q_5-\pi/2$ & 0               & $\pm80^\circ$\\
6 & Wrist yaw  & $-\pi/2$ & 0.0091 & $q_6-\pi/2$ & 0               & $\pm80^\circ$\\
\hline
\end{tabular}
\end{table}

\textit{Forward kinematics.} Multiplying the six transforms gives the pose of
the jaws for any set of joint values:
\begin{equation}
  \mathbf T(\mathbf q)={}^{0}\mathbf T_1\,{}^{1}\mathbf T_2\cdots{}^{5}\mathbf T_6
  =\begin{bmatrix}\mathbf R(\mathbf q)&\mathbf p(\mathbf q)\\\mathbf 0^\top&1\end{bmatrix}.
\end{equation}
The jaws point along $\hat{\mathbf a}(\mathbf q)$, the second column of
$\mathbf R(\mathbf q)$, and the tip is one jaw length
($\ell=\SI{10.2}{mm}$) further along:
\begin{equation}
  \mathbf p_{\mathrm{tip}}(\mathbf q)=\mathbf p(\mathbf q)+\ell\,\hat{\mathbf a}(\mathbf q).
\end{equation}
In the code, the same chain is written in the equivalent
product-of-exponentials form~\cite{lynch2017modern}, which makes the
derivatives below cheap to compute. The two forms agree to within
$10^{-15}$.

\textit{The RCM constraint.} Because yaw and pitch both rotate about the RCM
and insertion slides along a line through it, the end of the shaft (the
wrist) is always
\begin{equation}
  \mathbf p_w(\mathbf q)=(q_3-\SI{15.6}{mm})\,\hat{\mathbf s}(q_1,q_2),
\end{equation}
where $\hat{\mathbf s}$ is the direction the shaft points. In other words,
the shaft always passes through the RCM. Only the wrist and roll can angle
the jaws away from that line, and the wrist bends at most $80^\circ$. This is
what limits the directions from which a target can be approached.

\textit{Jacobian.} The Jacobian $\mathbf J(\mathbf q)$ is a $6\times6$ matrix
that says how fast the tip moves and turns for a small change in each joint.
It is what the solver below uses to decide which way to move the joints.

\subsection{Solving for the Joint Angles}\label{sec:ik}
Inverse kinematics (IK) is the reverse problem: given a target
$(\mathbf p,\mathbf a)$, find joint values $\mathbf q$ within the joint
limits such that
\begin{equation}
  \mathbf p_{\mathrm{tip}}(\mathbf q)=\mathbf p,
  \qquad
  \angle\big(\hat{\mathbf a}(\mathbf q),\mathbf a\big)\le\gamma .
\end{equation}
The jaws do not have to point exactly along $\mathbf a$, only within a cone
of half-angle $\gamma=20^\circ$ around it. Spinning the jaws about their own
axis is allowed freely, since that does not matter when pushing a stiff
segment straight through a hole.

The solver starts from a guess in which the shaft is aimed at the target with
the wrist straight, then improves it step by step. At each step it measures
two errors: how far the tip is from the target,
$\mathbf e_p=\mathbf p-\mathbf p_{\mathrm{tip}}$, and how far the jaw axis
is from the nearest direction inside the cone, $\mathbf e_\omega$. Stacking
these into one error vector $\mathbf e$, the joints are updated with the
damped least-squares step~\cite{wampler1986dls}
\begin{equation}
  \Delta\mathbf q=\big(\mathbf J^\top\mathbf J+\lambda\mathbf I\big)^{-1}\mathbf J^\top\mathbf e,
  \qquad \lambda=10^{-4}.
\end{equation}
This is the joint change that best reduces the error. The small $\lambda$
keeps the step from blowing up when the arm is near a configuration where
some motion is impossible. After each step, any joint pushed past its limit
is clamped back to the limit. One run uses 60 steps. It is repeated from four
different roll angles, and the best result is kept.

A target counts as \emph{reached} if the tip ends up within \SI{0.5}{mm}
of it and the jaws point inside the cone. As a check, the solver was
asked to reach 2000 random poses that the arm was known to be able to reach.
It found 99.7--99.9\% of them, so if anything it slightly under-counts
reachability.

\subsection{Scoring a Placement}\label{sec:score}
To rate a placement $\mathbf x$, every target is first rewritten relative to
the RCM:
\begin{equation}
  \mathbf p'_k=\mathbf R(\mathbf x)^\top(\mathbf p_k-\mathbf t),
  \qquad
  \mathbf a'_k=\mathbf R(\mathbf x)^\top\mathbf a_k .
\end{equation}
The IK is then solved for each target. To leave a safety margin, every
rotary joint is kept \SI{0.05}{rad} inside its limit, and insertion is
kept at least \SI{40}{mm}. A target counts as reachable, $s_k=1$, only if
the IK succeeds and all of the following hold:
\begin{itemize}
  \item the shaft (a \SI{4.2}{mm} radius cylinder from the RCM to the
    wrist) does not pass through either plate;
  \item the wrist and jaws (\SI{2.5}{mm} radius) do not pass through a
    plate, ignoring the last \SI{3}{mm} at the tip, which is meant to touch
    the target;
  \item no part of the tool goes below the table;
  \item the RCM is at least \SI{2}{cm} from the bracket and the table.
\end{itemize}
Otherwise $s_k=0$.

The placement's score adds up what fraction of corner points and holes are
reachable:
\begin{equation}\label{eq:objective}
  J(\mathbf x)=
  2\cdot\frac{\#\text{corner points reached}}{13}
  +\frac{\#\text{holes reached}}{421}
  +0.02\,\bar m .
\end{equation}
The corner is weighted twice because it is essential to the task. The last
term is a small bonus for keeping joints away from their limits. Here
$\bar m$ is the average, over reached targets, of how far the closest joint
is from a limit, as a fraction of its range: 0 at a limit, 0.5 at the middle.
Reaching everything gives a score of just over 3.

\subsection{Searching for the Best Placement}\label{sec:search}
Because the score counts targets, it jumps in steps as the placement moves.
Methods that follow a slope therefore do not work, so the group used
differential evolution~\cite{storn1997de} from SciPy~\cite{virtanen2020scipy}.
It keeps a population of 120 candidate placements spread over the search
range: the RCM within \SI{30}{cm} of the bracket and above the table, any
rotation about the vertical, and tilts up to a limit $\Theta$. Each
generation, every candidate is challenged by a new one, built by taking the
current best and moving it along the difference between two other random
candidates:
\begin{equation}
  \mathbf v=\mathbf x_{\mathrm{best}}+F\,(\mathbf x_{r_1}-\mathbf x_{r_2}),
  \qquad F\in[0.5,1).
\end{equation}
Some of the new candidate's six numbers are swapped back with the original
candidate's, and whichever of the two scores higher survives. Over about 60
generations, the population closes in on the best region.

For speed, the search scores candidates against every third hole plus all
corner points. The winner is then re-scored on all 434 targets.

\subsection{Results}\label{sec:results}
Two searches were run: placement~A with the tilt limited to
$\Theta=30^\circ$, and placement~B allowing up to $70^\circ$
(Table~\ref{tab:placement}). Both reach all 13 corner points and all 421
holes. Every joint solution for placement~A also fits within the slightly
narrower joint limits of the AMBF simulator's PSM model
(Table~\ref{tab:joints}), so the result can be tested directly in
simulation.

\begin{table}[t]
\centering
\caption{Best RCM placements found, in the bracket frame.}
\label{tab:placement}
\footnotesize
\begin{tabular}{@{}lcc@{}}
\hline
 & A ($\Theta=30^\circ$) & B ($\Theta=70^\circ$)\\
\hline
RCM position (mm)           & $(82.6,\,-87.3,\,127.8)$ & $(117.7,\,-81.1,\,76.1)$\\
$\psi,\theta_y,\theta_x$ ($^\circ$) & $-40.7,\ 27.1,\ 5.1$ & $-126.2,\ -23.4,\ -69.5$\\
Tool tilt from vertical     & $27.6^\circ$ & $71.2^\circ$\\
Corner / holes reached      & 13/13, 421/421 & 13/13, 421/421\\
Joint margin (mean / min)   & 0.314 / 0.167 & 0.331 / 0.167\\
Score $J$                   & 3.0063 & 3.0066\\
\hline
\end{tabular}
\end{table}

\begin{table}[t]
\centering
\caption{Joint ranges used by placement A across all 434 targets, compared
with the dVRK and AMBF limits (rad; insertion in m).}
\label{tab:joints}
\footnotesize
\setlength{\tabcolsep}{4pt}
\begin{tabular}{@{}lccc@{}}
\hline
Joint & Used (A) & dVRK & AMBF\\
\hline
Yaw $q_1$         & $[-0.654,\,-0.140]$ & $\pm1.588$ & $[-1.605,\,1.599]$\\
Pitch $q_2$       & $[-0.320,\,0.313]$  & $\pm0.925$ & $[-0.611,\,0.768]$\\
Insertion $q_3$   & $[0.124,\,0.186]$   & $[0,\,0.24]$ & $[0,\,0.24]$\\
Roll $q_4$        & $[-1.588,\,2.993]$  & $\pm4.538$ & $\pm3.054$\\
Wrist pitch $q_5$ & $[-0.618,\,0.610]$  & $\pm1.396$ & $\pm1.571$\\
Wrist yaw $q_6$   & $[-0.567,\,0.562]$  & $\pm1.396$ & $\pm1.571$\\
\hline
\end{tabular}
\end{table}

\begin{figure}[H]
    \centering
    \includegraphics[width=1\linewidth]{figures/coverage.png}
    \caption{Reachability of placement A. Left and centre: hole maps of the
    two plates, with reached targets in green; squares are the inside-corner
    points. Right: the plates, the RCM and a sample of tool poses (blue:
    shaft; orange: wrist and jaws).}
\label{fig:coverage}
\end{figure}

\textit{How strict can the alignment be?} Table~\ref{tab:cone} repeats the
check with a narrower cone. Placement~A still reaches everything when the
jaws must be within $5^\circ$ of the hole axis. Even at exact alignment it
reaches 12 of 13 corner points and 412 of 421 holes; the shaft clips the
other plate on the rest. Placement~B starts failing at $5^\circ$, because its
nearly horizontal shaft runs too close to the table.

\begin{table}[t]
\centering
\caption{Targets reached (corner / holes) as the allowed jaw misalignment
$\gamma$ shrinks.}
\label{tab:cone}
\footnotesize
\begin{tabular}{@{}ccc@{}}
\hline
$\gamma$ & A ($\Theta=30^\circ$) & B ($\Theta=70^\circ$)\\
\hline
$20^\circ$ & 13 / 421 & 13 / 421\\
$10^\circ$ & 13 / 421 & 13 / 421\\
$5^\circ$  & 13 / 421 & 13 / 419\\
$2^\circ$  & 12 / 420 & 13 / 416\\
$0^\circ$  & 12 / 412 & 12 / 409\\
\hline
\end{tabular}
\end{table}

\textit{Choice of placement.} The two placements score almost the same.
Placement~B, however, tips the arm $71^\circ$ over with the RCM only
\SI{76}{mm} above the table. The rest of the arm's body would have to fit in
that space, and this model does not check it for collisions. Placement~A
was therefore chosen for the first PSM. It also copes better with stricter
alignment, and it keeps every joint at least 16.7\% of its range away from
its limit.
